\documentclass[12pt]{article}
\usepackage[utf8]{inputenc}
\usepackage{geometry}
\geometry{a4paper, margin=1in}
\usepackage{hyperref}
\usepackage{enumitem}
\usepackage{titlesec}
\usepackage{xcolor}
\usepackage{fancyhdr}

% Header and Footer configuration
\pagestyle{fancy}
\fancyhf{}
\fancyhead[L]{\textbf{Week 04 Implementation Report} --- Palak Agarwal}
\fancyhead[R]{Page \thepage}
\fancyfoot[C]{\thepage}

\title{\textbf{Week 04 Extensive Implementation Report}}
\author{Palak Agarwal}
\date{Module: Student Information System (SIS)}

\begin{document}

\maketitle
\thispagestyle{fancy} % Apply header/footer to the title page

\section*{Objective}
Develop the complete Student Management foundation including database schema, backend APIs, basic CRUD operations, and comprehensive Student Lifecycle Workflows.

\vspace{0.5cm}
\hrule
\vspace{0.5cm}

% ------------------------- DAY 01 -------------------------
\section*{Day 01 --- Student Status Workflow (21st September 2026)}

\subsection*{1. What you have done?}
\begin{itemize}
    \item \textbf{Status Tracking}: Enhanced the \texttt{StudentStatus} enumeration to include \texttt{DROPPED} and \texttt{TRANSFERRED}.
    \item \textbf{History Logging}: Updated the \texttt{StudentStatusHistory} entity and \texttt{StudentStatusUpdateDto} to include an \texttt{effectiveDate} property for precise auditing.
    \item \textbf{System Identifiers}: Transitioned the Enrollment/Roll Number to be automatically generated by the system (\texttt{ENR-}) instead of relying on user input to avoid manual modification and improve data consistency.
\end{itemize}

\subsection*{2. What error you have encountered?}
\begin{itemize}
    \item \textbf{Constraint Validation}: Encountered exceptions when updating statuses due to missing effective dates in legacy API requests.
\end{itemize}

\subsection*{3. How did you resolve/overcome those error?}
\begin{itemize}
    \item \textbf{Resolution}: Introduced a fallback in the \texttt{StudentServiceImpl} to automatically default the \texttt{effectiveDate} to \texttt{LocalDate.now()} if not explicitly provided in the payload.
\end{itemize}

\subsection*{4. What did you learn?}
\begin{itemize}
    \item \textbf{Audit Completeness}: Emphasized the importance of effective dating in historical audit tables, ensuring timelines for academic changes remain accurate regardless of system update times.
\end{itemize}

\vspace{0.5cm}
\hrule
\vspace{0.5cm}

% ------------------------- DAY 02 -------------------------
\section*{Day 02 --- Student APIs \& Testing (22nd September 2026)}

\subsection*{1. What you have done?}
\begin{itemize}
    \item \textbf{API Finalization}: Stabilized the complete suite of Student CRUD APIs, integrating system-generated IDs and effective dating.
    \item \textbf{Test Coverage}: Wrote unit and integration tests across \texttt{StudentControllerTest} using MockMvc to validate endpoints like GET by Roll Number, Pagination, and Status Updates.
\end{itemize}

\subsection*{2. What error you have encountered?}
\begin{itemize}
    \item \textbf{Test Data Conflicts}: Tests were failing due to non-unique roll numbers being inserted into the test context.
\end{itemize}

\subsection*{3. How did you resolve/overcome those error?}
\begin{itemize}
    \item \textbf{Resolution}: Mocked the roll number generation logic in tests to guarantee unique identifiers and avoid unique constraint violations in H2 memory.
\end{itemize}

\subsection*{4. What did you learn?}
\begin{itemize}
    \item \textbf{Test Isolation}: Understood the necessity of properly cleaning up state and using unique identifiers between independent test cases.
\end{itemize}

\vspace{0.5cm}
\hrule
\vspace{0.5cm}

% ------------------------- DAY 03 -------------------------
\section*{Day 03 --- Student Enrollment (23rd September 2026)}

\subsection*{1. What you have done?}
\begin{itemize}
    \item \textbf{Schema Connection}: Engineered the \texttt{StudentEnrollment} entity to establish a linkage between a \texttt{Student} and the academic structure (\texttt{AcademicYear}, \texttt{Program}, \texttt{Batch}, \texttt{Semester}).
    \item \textbf{Features Added}: Incorporated fields for \texttt{enrollmentDate}, \texttt{enrollmentNumber}, and an isolated \texttt{status} to manage re-enrollments and transfers.
\end{itemize}

\subsection*{2. What error you have encountered?}
\begin{itemize}
    \item \textbf{Duplicate Enrollments}: Faced challenges designing database constraints to prevent duplicate enrollments for the same student in the same semester.
\end{itemize}

\subsection*{3. How did you resolve/overcome those error?}
\begin{itemize}
    \item \textbf{Resolution}: Added a \texttt{@UniqueConstraint} at the entity level on \texttt{student\_id}, \texttt{academic\_year\_id}, and \texttt{semester\_id} to strictly enforce data integrity in the database.
\end{itemize}

\subsection*{4. What did you learn?}
\begin{itemize}
    \item \textbf{Relational Mapping}: Learned how to structurally relate broad institutional models to individual profiles, utilizing multi-column unique indexes to prevent logic faults.
\end{itemize}

\vspace{0.5cm}
\hrule
\vspace{0.5cm}

% ------------------------- DAY 04 -------------------------
\section*{Day 04 --- Student Academic Records (24th September 2026)}

\subsection*{1. What you have done?}
\begin{itemize}
    \item \textbf{Academic History}: Designed the \texttt{StudentAcademicRecord} entity to maintain semester-wise academic history (GPA, CGPA).
    \item \textbf{Integration}: Connected the academic performance metrics back to the core student profile.
\end{itemize}

\subsection*{2. What error you have encountered?}
\begin{itemize}
    \item \textbf{Module Overlap}: Difficulty separating the concern of "Marks/Results" from "Student Records," creating a risk of data duplication.
\end{itemize}

\subsection*{3. How did you resolve/overcome those error?}
\begin{itemize}
    \item \textbf{Resolution}: Resolved by referencing high-level academic performance metrics (GPA/CGPA) directly, ensuring the entity acts as a consumer of the examination data rather than a source of truth for individual subject marks.
\end{itemize}

\subsection*{4. What did you learn?}
\begin{itemize}
    \item \textbf{Domain-Driven Design (DDD)}: Discovered the importance of Domain-Driven Design principles—specifically avoiding the duplication of contexts across independent modules.
\end{itemize}

\vspace{0.5cm}
\hrule
\vspace{0.5cm}

% ------------------------- DAY 05 -------------------------
\section*{Day 05 --- Student Category \& Classification (25th September 2026)}

\subsection*{1. What you have done?}
\begin{itemize}
    \item \textbf{Classification Schema}: Implemented the \texttt{StudentCategory} entity to classify students (e.g., General, OBC, SC, ST) dynamically.
\end{itemize}

\subsection*{2. What error you have encountered?}
\begin{itemize}
    \item \textbf{Hardcoding Limitations}: Initially planned to use a simple Java Enum for categories, which limits future expandability if the institution wants to add a new scholarship or minority classification.
\end{itemize}

\subsection*{3. How did you resolve/overcome those error?}
\begin{itemize}
    \item \textbf{Resolution}: Transitioned from an Enum structure to a database-driven dictionary table (\texttt{student\_categories}) allowing dynamic administrative management without code deployments.
\end{itemize}

\subsection*{4. What did you learn?}
\begin{itemize}
    \item \textbf{Dynamic Configuration}: Mastered the principle of dynamic configurations versus static logic, realizing that institutional parameters change frequently and should be table-driven.
\end{itemize}

\vspace{0.5cm}
\hrule
\vspace{0.5cm}

% ------------------------- DAY 06 -------------------------
\section*{Day 06 --- Student Lifecycle Management (26th September 2026)}

\subsection*{1. What you have done?}
\begin{itemize}
    \item \textbf{Lifecycle Workflow}: Implemented a rigorous state machine for student statuses within \texttt{StudentServiceImpl}.
    \item \textbf{Validation Logic}: Created \texttt{validateStatusTransition} to enforce proper state changes (e.g., preventing \texttt{ALUMNI} from changing to \texttt{DROPPED} or ensuring \texttt{GRADUATED} strictly flows to \texttt{ALUMNI}).
\end{itemize}

\subsection*{2. What error you have encountered?}
\begin{itemize}
    \item \textbf{Invalid State Change}: An initial test allowed terminal states like \texttt{ALUMNI} and \texttt{DROPPED} to be reactivated, violating business rules.
\end{itemize}

\subsection*{3. How did you resolve/overcome those error?}
\begin{itemize}
    \item \textbf{Resolution}: Enforced terminal state checks within the \texttt{switch} expression of the transition validation logic, returning \texttt{false} for terminal state triggers and throwing a \texttt{BusinessException}.
\end{itemize}

\subsection*{4. What did you learn?}
\begin{itemize}
    \item \textbf{Domain Driven Rules}: Learned that enforcing strict lifecycle transition matrices in the backend shields the database from logically inconsistent data states.
\end{itemize}

\end{document}
